\section{Design Thinking: Mapa de Empat\'ia y definici\'on del problema}

\subsection{Prop\'osito del lienzo}
Tomando el \emph{design thinking} como base para la ingenier\'ia de requisitos \cite{ireb2026}, en esta secci\'on se arma el Mapa de Empat\'ia del \textbf{Cajero / Vendedor de Mostrador (SH-01)}. Se usa para ordenar lo que salga del ensayo de entrevista con ese rol: qu\'e dice, qu\'e hace, qu\'e piensa y qu\'e siente cuando atiende en el mostrador. Las celdas se llenaron con lo que se dijo en el ensayo \texttt{GA-TL-02-ENS-01}, cuyo audio est\'a registrado en \texttt{GA/TL-02/evidencias/MANIFIESTO.md}. Las frases entre comillas son textuales.

\subsection{Mapa de Empat\'ia (Perfil: SH-01 --- Cajero / Vendedor de Mostrador)}

\begin{tcolorbox}[colback=boxbg,colframe=midblue,boxrule=0.8pt,arc=2mm,title={\textbf{Perfil analizado: Cajero / Vendedor de Mostrador (SH-01) --- Ferreter\'ia el\'ectrica en Quevedo}}]
\textbf{Escenario ensayado:} Atenci\'on en mostrador durante horas pico (ma\~nana y mediod\'ia), b\'usqueda de materiales solicitados con jerga local, asesor\'ia t\'ecnica sobre reemplazos y cobro con facturaci\'on electr\'onica.\\
\textbf{Fuente de respaldo:} Audio del ensayo \texttt{GA-TL-02-ENS-01} registrado en \texttt{GA/TL-02/evidencias/MANIFIESTO.md}.
\end{tcolorbox}

\vspace{4pt}
{\small
\noindent\begin{tabularx}{\linewidth}{|Y|Y|}
\hline
\thdr{¿Qu\'e DICE? (Expresiones del ensayo)} & \thdr{¿Qu\'e HACE? (Acciones en mostrador)}\\
\hline
\begin{itemize}[leftmargin=12pt,itemsep=2pt,topsep=2pt]
    \item ``La gente llega pidiendo con nombres de calle como \emph{taipe}, \emph{breaker} o \emph{boquilla}, y a veces toca adivinar el amperaje.'' \textit{(ENS-01)}
    \item ``Cuando se hace fila al mediod\'ia, perder tres minutos buscando un reemplazo en bodega o preguntando al due\~no molesta al cliente.'' \textit{(ENS-01)}
    \item ``Si el internet se cae o el SRI se pone lento, toca anotar en libreta y cuadrar despu\'es.'' \textit{(ENS-01)}
\end{itemize}
&
\begin{itemize}[leftmargin=12pt,itemsep=2pt,topsep=2pt]
    \item Alterna constantemente entre cobrar en caja, revisar existencias en bodega y responder dudas t\'ecnicas de clientes particulares.
    \item Consulta al administrador, revisa cat\'alogos impresos o busca en el tel\'efono cuando desconoce una equivalencia el\'ectrica.
    \item Prioriza el uso de teclado y lector de barras por el espacio reducido en el mostrador.
\end{itemize}\\
\hline
\thdr{¿Qu\'e PIENSA Y SIENTE? (Emociones y temores)} & \thdr{¿Qu\'e VE y OYE? (Entorno f\'isico y social)}\\
\hline
\begin{itemize}[leftmargin=12pt,itemsep=2pt,topsep=2pt]
    \item Siente presi\'on cuando se acumulan varios clientes apurados en la fila y debe verificar datos t\'ecnicos al mismo tiempo.
    \item Teme recomendar un material incompatible que provoque una devoluci\'on, un da\~no en la instalaci\'on del cliente o un reclamo del administrador.
    \item Percibe que un vendedor nuevo tarda meses en memorizar equivalencias y nombres locales del cat\'alogo.
\end{itemize}
&
\begin{itemize}[leftmargin=12pt,itemsep=2pt,topsep=2pt]
    \item Observa un mostrador estrecho, cajas de materiales el\'ectricos y clientes revisando la hora mientras esperan su factura.
    \item Escucha pedidos apurados de maestros electricistas que exigen precisi\'on y clientes particulares que no saben el nombre t\'ecnico del producto.
    \item Oye quejas cuando un producto aparece con saldo en pantalla pero no existe f\'isicamente en percha.
\end{itemize}\\
\hline
\thdr{DOLORES / FRUSTRACIONES} & \thdr{RESULTADOS / NECESIDADES}\\
\hline
\begin{itemize}[leftmargin=12pt,itemsep=2pt,topsep=2pt]
    \item Demoras de varios minutos para confirmar compatibilidad de sustitutos cuando un art\'iculo est\'a agotado.
    \item Desfases entre el inventario registrado y la disponibilidad real en bodega.
    \item Interrupciones operativas cuando falla la conectividad o el servicio tributario del SRI.
\end{itemize}
&
\begin{itemize}[leftmargin=12pt,itemsep=2pt,topsep=2pt]
    \item Responder con rapidez y seguridad sobre alternativas compatibles sin depender siempre del vendedor m\'as antiguo.
    \item Confiar en que el stock consultado corresponde a la existencia f\'isica real.
    \item Mantener la continuidad de la atenci\'on en caja incluso ante fallos externos de red.
\end{itemize}\\
\hline
\end{tabularx}
}

\subsection{Definici\'on del problema centrada en el usuario}
La definici\'on del problema se redacta estrictamente desde la necesidad operativa del trabajador en mostrador, evitando imponer una soluci\'on tecnol\'ogica espec\'ifica:

\begin{tcolorbox}[colback=white,colframe=uteqgreen,boxrule=0.8pt,arc=1.5mm]
\textbf{Enunciado de punto de vista (POV):}\\
El cajero y vendedor de mostrador de un establecimiento de suministros el\'ectricos en Quevedo necesita identificar en segundos productos compatibles, disponibles en bodega y reconocibles por su nombre coloquial local, porque durante las horas pico las consultas t\'ecnicas manuales y los desfases de inventario retrasan la atenci\'on y generan riesgo de vender materiales incompatibles.\\[8pt]
\textbf{Pregunta generadora («¿C\'omo podr\'iamos...?»):}\\
\emph{¿C\'omo podr\'iamos ayudar al cajero y vendedor de mostrador a identificar con rapidez y certeza alternativas el\'ectricas compatibles y disponibles —incluso cuando el cliente usa nombres coloquiales— sin detener el flujo de atenci\'on en horas pico ni depender de consultas manuales prolongadas?}
\end{tcolorbox}
